% !TeX root = ../main.tex
% Fonti: Lezione 9, pp. 1--4; Lezione 10, pp. 1--5.
\chapter{Equazioni, prodotti e quozienti}
\fonti{Lezioni 9 e 10}

\section{Equazioni parametriche e cartesiane}
Una base $u_1,\ldots,u_r$ di un sottospazio $U\subseteq K^m$ dà una
descrizione \emph{parametrica}:
\[
U=\{a_1u_1+\cdots+a_ru_r:a_1,\ldots,a_r\in K\}.
\]
Una descrizione \emph{cartesiana} specifica invece le coordinate
mediante equazioni lineari omogenee. Per passare dalla prima alla
seconda si eliminano i parametri; per il verso inverso si scelgono
coordinate libere e si risolvono le equazioni restanti. La
omogeneità è essenziale: l'insieme delle soluzioni contiene il
vettore nullo ed è chiuso per le combinazioni lineari.

\begin{esercizio}
Riprendiamo
$U=\langle(1,-1,0,2),(2,0,-1,3)\rangle\subseteq\R^4$.
Scrivendo
\[
(x_1,x_2,x_3,x_4)
=a(1,-1,0,2)+b(2,0,-1,3)
\]
si hanno $x_1=a+2b$, $x_2=-a$, $x_3=-b$ e $x_4=2a+3b$.
Le coordinate seconda e terza danno $a=-x_2$, $b=-x_3$; sostituendo
nelle altre due otteniamo
\[
U=\{x\in\R^4:x_1+x_2+2x_3=0,\quad
2x_2+3x_3+x_4=0\}.
\]
Le condizioni sono anche sufficienti: se $x$ le soddisfa, i valori
$a=-x_2$ e $b=-x_3$ ricostruiscono tutte e quattro le coordinate,
quindi $x\in U$.
\end{esercizio}

\section{Prodotto di spazi vettoriali}
\begin{proposizione}
Se $V$ e $W$ sono spazi vettoriali sullo stesso campo $K$, il prodotto
cartesiano $V\times W$ è uno spazio vettoriale con
\[
(v_1,w_1)+(v_2,w_2)=(v_1+v_2,w_1+w_2),\qquad
\lambda(v,w)=(\lambda v,\lambda w).
\]
Se $V$ e $W$ hanno dimensioni finite, allora
$\dim(V\times W)=\dim V+\dim W$.
\end{proposizione}
\begin{proof}
Il vettore nullo è $(0_V,0_W)$ e l'opposto di $(v,w)$ è
$(-v,-w)$. Associatività, commutatività e distributività seguono
verificando separatamente le due componenti. Per la dimensione,
siano $v_1,\ldots,v_r$ una base di $V$ e $w_1,\ldots,w_s$ una base
di $W$. La famiglia
\[
(v_1,0_W),\ldots,(v_r,0_W),
(0_V,w_1),\ldots,(0_V,w_s)
\]
genera il prodotto, perché $(v,w)$ si scompone nelle due espansioni
in base. È indipendente: una combinazione nulla produce nella prima
componente una relazione tra i $v_i$ e nella seconda una relazione
tra i $w_j$; tutti i coefficienti sono dunque nulli. La base così
costruita ha $r+s$ elementi.
\end{proof}

\section{Classi laterali e spazio quoziente}
\begin{definizione}
Sia $W$ un sottospazio di $V$. Due vettori sono equivalenti modulo
$W$ se $v_1-v_2\in W$. La classe di $v$ è
\[
[v]=v+W=\{v+w:w\in W\}.
\]
L'insieme delle classi è scritto $V/W$ e si chiama
\emph{spazio quoziente}.
\end{definizione}
La relazione è riflessiva perché $v-v=0_V\in W$; è simmetrica perché
$-(v_1-v_2)\in W$; è transitiva perché la somma di
$v_1-v_2$ e $v_2-v_3$ appartiene a $W$. Due classi coincidono
esattamente quando i rispettivi rappresentanti differiscono di un
elemento di $W$.

Definiamo $[v_1]+[v_2]=[v_1+v_2]$ e
$\lambda[v]=[\lambda v]$. Occorre controllare che il risultato non
dipenda dai rappresentanti scelti. Se $v_i'\in[v_i]$, allora
$v_i'-v_i\in W$, per $i=1,2$, e
\[
(v_1'+v_2')-(v_1+v_2)
=(v_1'-v_1)+(v_2'-v_2)\in W.
\]
Inoltre $\lambda v'-\lambda v=\lambda(v'-v)\in W$.
Le operazioni sono quindi ben definite; gli assiomi di spazio
vettoriale passano da $V$ alle classi. Il vettore nullo del
quoziente è $[0_V]=W$, non il solo vettore $0_V$.

\begin{esempio}
Se $W$ è l'asse orizzontale di $\R^2$, due punti $(x,y)$ e
$(x',y')$ sono equivalenti se $y=y'$. Ogni classe è una retta
orizzontale: il quoziente registra solo la coordinata verticale.
\end{esempio}
\begin{schema}
\begin{tikzpicture}[scale=.7]
  \draw[asse] (-3,0)--(3.3,0) node[right] {$W$};
  \draw[asse] (0,-.8)--(0,2.7);
  \draw[spazio] (-2.6,1.45)--(2.7,1.45);
  \fill[coloreschema] (-1.45,1.45) circle (2pt)
      node[above] {$v$};
  \fill[coloreschema] (1.25,1.45) circle (2pt)
      node[above] {$v'$};
  \node[right] at (2.7,1.45) {$[v]=[v']$};
\end{tikzpicture}
\end{schema}

\begin{teorema}
Se $V$ è di dimensione finita e $W\subseteq V$ è un sottospazio,
allora $\dim(V/W)=\dim V-\dim W$.
\end{teorema}
\begin{proof}
Scegliamo una base $w_1,\ldots,w_r$ di $W$ e completiamola
con $v_1,\ldots,v_s$ fino a una base di $V$. Le classi
$[v_1],\ldots,[v_s]$ generano $V/W$: ogni $v\in V$ si scrive
$v=\sum_i a_iw_i+\sum_j b_jv_j$, e nel quoziente la prima somma
ha classe nulla. Se $\sum_j b_j[v_j]=[0_V]$, allora
$\sum_j b_jv_j\in W$, quindi
$\sum_j b_jv_j=\sum_i a_iw_i$. L'indipendenza della base di $V$
impone $a_i=b_j=0$ per tutti gli indici. Le $s$ classi sono dunque
una base; poiché $\dim V=r+s$ e $\dim W=r$, segue la formula.
\end{proof}
